% §2.6 — NOₓ Reduction Input-Output Model: σ Elimination and Parametrization
% Sources: Thesis/secs/4-NOx_mdl/1_NOx_r_mdl.tex
%          Thesis/secs/4-NOx_mdl/1-subs/change_in_nox_model.tex
%          Thesis/secs/4-NOx_mdl/1-subs/parametrization.tex
%          Thesis/secs/4-NOx_mdl/1-subs/ident_form.tex
%          Thesis/secs/4-NOx_mdl/25_switching_condition.tex
% ===========================================================
\subsection{$\nox$ Reduction Input-Output Model \label{sec::nox_io_model}}

% --- Control volume figure ---
% \begin{figure}[!ht]
%     \centering
%     \includegraphics[width=\figWidth]{./figs/2-modeling/ctrl_vol.png}
%     \caption{SCR-ASC system abstraction}
%     \label{fig::ctrl_vol}
% \end{figure}

The internal state $\sigma(k)$ is not directly measurable. To obtain an identifiable input-output model, $\sigma$ is
eliminated using the $\nox$ dynamics. Defining:
\begin{align}
    \eta(k) &= u_1(k-1) - x_1(k)
    \label{eqn::eta_def}
\end{align}
where $u_1(k) = \con{\nox}^{in}(k)$ and $x_1(k) = \con{\nox}^{out}(k)$, the quantity $\eta(k)$ represents the change
in $\nox$ concentration due to catalytic reduction.

% --- σ Elimination ---
\subsubsection{$\sigma$ Elimination}
From equation~(\ref{eqn::NOX_process}), the average surface concentration can be expressed as:
\begin{align}
    \sigma(k) &= \frac{\eta(k+1)}{\tau(k) k_{s2v} k_{scr}(k) u_1(k)}
    \label{eqn::sigma_elim}
\end{align}
Substituting into the adsorption dynamics~(\ref{eqn::sigma_process}) and defining the lumped
oxidation-desorption rate constant $k_{od} = k_{oxi} + k_{des}$, we obtain the recursive $\eta$-dynamics:
% ===
% TODO: Fill in the full η(k+1) equation from thesis eqn (4.8)
% ===
\begin{align}
    \eta(k+1) =& \eta(k) \lrb{\frac{u_1(k)}{F(k)}} \lrb{\frac{F(k-1)}{u_1(k-1)}} \notag \\
        &- \eta(k) \lrb{\frac{u_1(k)}{F(k)}} \lrb{\frac{F(k-1)}{u_1(k-1)}}
            t_s k_{ads}(k-1) \con{\ammonia}^{in}(k-1) \notag \\
        &- \eta(k) \lrb{\frac{u_1(k)}{F(k)}} \lrb{\frac{F(k-1)}{u_1(k-1)}}
            t_s k_{scr}(k-1) u_1(k-1) \notag \\
        &- \eta(k) \lrb{\frac{u_1(k)}{F(k)}} \lrb{\frac{F(k-1)}{u_1(k-1)}}
            t_s k_{od}(k-1) \notag \\
        &+ t_s k_{s2v} \times \lrf{\Gamma(k-1) \tau(k) u_1(k) \con{\ammonia}^{in}(k-1)}
            \times \lr{k_{scr}(k) k_{ads}(k-1)}
    \label{eqn::disc_NOx}
\end{align}

% --- Parametrization ---
\subsubsection{Parametrization with Chebyshev Temperature Models}

Defining the \textit{flow-scaled fractional $\nox$ reduction}:
\begin{align}
    \bar \eta_F (k) &= \frac{\eta(k)}{u_{1F}(k-1)} = \frac{F(k-1)}{u_1(k-1)} \eta(k)
    \label{eqn::eta_bar_def}
\end{align}
where $u_{1F}(k) = u_1(k)/F(k)$, and applying the Chebyshev basis polynomials
(equations~\ref{eqn::phi_1}--\ref{eqn::phi_2}) to parametrize the rate constants, the $\eta$-dynamics can be written in
compact linear-in-parameters form:
\begin{align}
    \bar \eta_F (k+1) = \bar \eta_F (k) &+ u_{2F}(k-1)\phi_2(k-1) \theta_\Gamma
        \label{eqn::t2}\\
    &- u_{2F}(k-1) \bar \eta_F (k) \phi_1(k-1) \theta_{\eta_{ads}}
        \label{eqn::t3}\\
    &- \bar \eta_F (k) \phi_1(k-1) \theta_{\eta_{od}}
        \label{eqn::t4}\\
    &- u_1(k-1) \bar \eta_F (k) \phi_1(k-1) \theta_{\eta_{scr}}
        \label{eqn::t5}
\end{align}
where $u_{2F}(k) = u_2(k)/F(k)$, and $u_2(k) = u_{inj}(k)$ is the urea injection rate. The four parameter groups are:
\begin{inparaenum}[(1)]
    \item $\theta_\Gamma$: adsorption capacity (product of $k_{scr}$, $k_{ads}$, and $\Gamma$);
    \item $\theta_{\eta_{ads}}$: adsorption rate;
    \item $\theta_{\eta_{od}}$: combined oxidation-desorption rate; and
    \item $\theta_{\eta_{scr}}$: SCR reaction rate.
\end{inparaenum}

% --- Identifiable form ---
\subsubsection{Identifiable Form}
The model can be written compactly as:
\begin{align}
    \bar \eta_F (k+1) &= \bar \eta_F (k) + \phi_{NO_x}^T(k) \theta_{NO_x}
    \label{eqn::eta_compact}
\end{align}
where $\phi_{NO_x}(k)$ is the regressor vector that depends only on measured quantities (inlet and outlet $\nox$
concentrations, exhaust flow rate, temperature, and urea injection rate), and $\theta_{NO_x} = \lrb{\theta_{\eta_{ads}}^T
\; \theta_{\eta_{od}}^T \; \theta_{\eta_{scr}}^T \; \theta_\Gamma^T}^T$ is the parameter vector. This linear-in-parameters
structure eliminates the sum-of-exponentials identifiability issue found in continuous-time models.
